\section{Introduction}
\label{sec:intro}

Modern football already instruments one of the two bodies in every contact.
The match ball used at the 2022 World Cup carries an inertial sensor suspended
at its centre, sampled at 500\,Hz, whose sudden change of acceleration marks
the instant the ball is touched; that timestamp, fused with limb tracking from
cameras, is what made semi-automated offside possible and what settled whether
a striker's head had grazed a goal-bound cross \cite{fifa2022connected}. The
ball, in other words, knows \emph{when} it was hit. It does not know
\emph{what} hit it, nor whether the boot that hit it had struck an opponent a
few milliseconds earlier --- which is precisely the question a referee must
answer to call a foul.

This paper asks whether the other body in the contact can be instrumented in
the same spirit. A small IMU in the boot (in the heel counter or under the
insole) sees every impact the foot delivers. Three signatures matter: the
\emph{ball}, a short, hard, forward-axis deceleration (foot--ball contact lasts
about ten milliseconds \cite{nunome2006impact}) followed by a high-frequency
ringing from the stiff boot--ball interface; the \emph{player}, a longer,
softer, oblique pulse as soft tissue compresses, followed by a low-frequency
push as the bodies keep moving into each other; and the \emph{ground}, a
periodic, mostly vertical foot strike that is the boot's normal background.
The interesting --- and contested --- case is the fourth: the boot hits ball and
opponent almost at once, and the call turns on which came first.

We treat this as a measurement problem with a decision-theoretic ending, and
we answer it by simulation. There is no public dataset of instrumented fouls,
so we build a physics-based generator for the boot's signal and pass it through
an explicit sensor model, using real ankle-IMU recordings of walking, running
and playing football \cite{reiss2012introducing} as an alternative background
to check that the synthetic gait is not doing the work. The study then runs in
four stages:

\begin{enumerate}
\item \textbf{Detection.} Can the boot tell ball from player from ground, and
  at what sampling rate and sensor grade? We compare the threshold logic of a
  connected ball (``a sudden change of the acceleration vector'') with a
  gradient-boosting classifier on interpretable pulse features.
\item \textbf{Timing and ordering.} How precisely is a contact onset timed,
  and --- combining the boot with the ball's own sensor --- how far apart must a
  ball contact and a player contact be for their order to be recovered?
\item \textbf{The anomalous overlap.} When both pulses fall inside one pulse
  width, a two-template non-negative least-squares fit separates them; we map
  where it works and where the physics says nothing can.
\item \textbf{The probabilistic call.} A knowledge base of situational base
  rates (how often a sliding tackle plays the ball first) is fused with the
  boot posterior, the ball-sensor timing and the separator's order into a
  calibrated probability that the ball was played first, and a cost-weighted
  rule with a video-review band decides.
\end{enumerate}
A fifth section places the contact in space: foot-mounted inertial navigation
with zero-velocity updates \cite{foxlin2005pedestrian} and sparse
ultra-wideband anchor fixes, evaluated as the probability of putting a contact
on the wrong side of the penalty-area line.

Our contribution is a quantified map of what a boot sensor can and cannot
settle, not a product. The numbers are specific to a simulator whose
assumptions we state and whose knobs we expose; every figure and every number
in the text regenerates from one command.
